\documentclass[12pt]{article}
\usepackage[utf8]{inputenc}
\usepackage[T1]{fontenc}
\usepackage{hyperref}
\title{Waamwang Language Resources for K-12 Curriculum Development}
\author{Compiled by Claude Sonnet 5.5 (claude-sonnet-5-5) for Andr\'as Kornai}
\date{}
\begin{document}
\maketitle

\section{Introduction}
Waamwang (also spelled Wamoang; ISO 639-3 wmn, Glottocode waam1236) is an extinct Austronesian (Oceanic) language of New Caledonia, historically tied to the Voh area. The ASJP database lists it as extinct since 1959. Searches found almost nothing beyond catalogue entries and a short wordlist. No native-language teaching materials were found.

\section{Primary Educational Resources}
No primers, textbooks or school materials for Waamwang were found. A curriculum would have to be built on reconstruction and on materials from neighbouring Kanak languages.

\section{Linguistic and Reference Materials}
\begin{itemize}
\item Wikipedia stub on the language, which cites Ethnologue: \url{https://en.wikipedia.org/wiki/Waamwang_language}
\item ASJP wordlist entry, drawn from Leenhardt's 1946 survey: \url{https://asjp.clld.org/languages/WAAMWANG}
\item Glottolog entry (waam1236): \url{https://glottolog.org/resource/languoid/id/waam1236.bigmap.html}
\end{itemize}
No grammar or dictionary was found in the search results.

\section{Religious and Community Resources}
None found.

\section{Digital and Audio Resources}
No audio, video or software resources were found.

\section{Curriculum Development Context}
Waamwang has no known speakers, and the only documentation found is a historical wordlist. A heritage curriculum could teach the language's history and place within New Caledonian Kanak heritage. Using the language itself would need further archival research, for example into Leenhardt's materials and work on neighbouring Voh languages such as Pwapw\^a. Neither of those was verified in these searches.

\end{document}